\paragraph{Area 2 – Dati dei docenti}~\\
\subsubsection{[ED2.1] Dati del docente}
Per ogni docente il sistema registra: \textbf{e-mail istituzionale}, \textbf{nome} e \textbf{cognome}. Tutti i dati sono obbligatori e non possono essere vuoti o composti solo da spazi.\\
Priorità: MUST\\Fonte: Traccia

\subsubsection{[ED2.2] E-mail come codice identificativo}
L'e-mail identifica il docente in modo univoco ed è il suo codice identificativo; non è modificabile dopo l'inserimento. Due e-mail che differiscono solo per maiuscole e minuscole sono considerate uguali; l'e-mail viene memorizzata in minuscolo.\\
Priorità: MUST\\Fonte: Traccia, C2, C10

\subsubsection{[ED2.3] Formato dell'e-mail}
L'e-mail deve avere la forma \texttt{parte-locale@unisa.it}, con una parte locale non vuota e priva di spazi. E-mail di altri domini (compreso \texttt{studenti.unisa.it}) vengono rifiutate.\\
Priorità: MUST\\Fonte: C15
